\documentclass[10pt,a4paper]{amsart}
\usepackage[margin=19mm]{geometry}
\usepackage{amsmath,amssymb,mathtools}
\usepackage[scr=rsfs,cal=euler]{mathalfa}
\usepackage{xfrac}
\usepackage[unicode,hidelinks,bookmarks=false]{hyperref}
\newcommand{\ie}{i.e.\ }
\newcommand{\cf}{cf.\ }
\newcommand{\resp}{respectively\ }
\newcommand{\Subsection}{\subsection}
\providecommand{\nobreakdash}{\nobreak\hbox{-}}
\setlength{\emergencystretch}{2em}
\multlinegap0pt
\usepackage{bm}
\usepackage{bbm}
\usepackage{dsfont}
\usepackage{xspace}
\usepackage{cancel}
\usepackage{mathtools}
\usepackage{amsthm}
\makeatletter
\@ifundefined{theorem}{\newtheorem{theorem}{\bf Theorem}}{}
\@ifundefined{lemma}{\newtheorem{lemma}[theorem]{\bf Lemma}}{}
\makeatother
\DeclareMathOperator{\ord}{ord}
\DeclareMathOperator{\supp}{supp}
\newcommand{\N}{\mathbb N}
\renewcommand{\t}{\, | \,}
\title[On a classical zero-sum invariant]{On a classical zero-sum invariant}
\author{Alfred Geroldinger, Wenkai Yang}
\date{Source: arXiv:2608.19090v1 (CC BY 4.0)}
\begin{document}
\maketitle

\noindent\emph{Source and licence.} On a classical zero-sum invariant. Version arXiv:2608.19090v1 (CC BY 4.0), distributed under the Creative Commons Attribution 4.0 International licence, as recorded in the arXiv licence metadata for this exact version and shown on the abstract page https://arxiv.org/abs/2608.19090v1. Full rights notice: licenses/arxiv-2608.19090-CC-BY-4.0.txt.\par\smallskip

\noindent\textit{Editorial scope.} Counted unit: the Theorem whose source label is \texttt{6.5} in the article (source0.tex lines 1653--1670, source label \texttt{6.5}), delivered together with the article's own complete original proof of it (source0.tex lines 1672--2038, one \texttt{proof} environment of 367 lines). No mathematics is written, repaired or re-derived: statement and proof are the article's own text line for line. Required context retained uncounted: 2.3 (lines 194--205); def-nu(S) (lines 1519--1521); 6.2 (lines 1525--1532). Every definition, display and label of those lines is retained unchanged. Omitted from this package and not redistributed: 1--193, 206--1518, 1522--1524, 1533--1652, 1671--1671, 2039--2370. Nothing inside a retained excerpt is omitted or altered: every retained body is delivered exactly as the article's lines are written, so no omission receipt of any kind accompanies this package. Citations: the retained excerpts carry no citation command at all, so no bibliography entry of the article is transcribed and no reference is redistributed. Reference adaptation: none was needed and none was made; every reference inside the delivered unit resolves inside it, so no unresolved reference or citation is produced. Printed slips kept literal: no printed word, symbol, hypothesis or number of the counted statement or of its proof was altered. Layout and class: the article's own class is \texttt{amsart} with packages including amsmath, amssymb, color, hyperref, xcolor; the wrapper is the lane's pinned \texttt{amsart} editorial wrapper at 10pt on A4, which restates the theorem-like environments the retained text needs and adds the article's own macro definitions that the retained text needs (\texttt{\textbackslash N}, \texttt{\textbackslash ord}, \texttt{\textbackslash supp}, \texttt{\textbackslash t}), with extra wrapper packages (bm, bbm, dsfont, xspace, cancel, mathtools). The delivered numbering is the unit's own. Assets: no retained span contains a figure environment, a graphics-inclusion command, a PDF-page inclusion, a table or a TikZ picture; no image file is copied into this package, the package declares no asset dependency and no image file is redistributed with it. Licence: version arXiv:2608.19090v1 of this article is distributed under the Creative Commons Attribution 4.0 International licence, as recorded in the arXiv licence metadata for this exact version and shown on the abstract page https://arxiv.org/abs/2608.19090v1. A full rights notice is delivered at licenses/arxiv-2608.19090-CC-BY-4.0.txt. Characters: every retained excerpt is pure ASCII, so the delivered engine, XeLaTeX with its default Latin Modern text font, reports no missing character.
\begin{lemma}\label{2.3}
Let $T \in \mathcal F (G)$ be zero-sum free.
\begin{enumerate}
\item  If  $|T|=\mathsf d (G)$, then $G^{\bullet} \setminus \Sigma (T)= \emptyset$ and $S = (-\sigma (T))T \in \mathcal A_{\max}(G)$.

\item If $|G^{\bullet} \setminus \Sigma (T)| \le 1$, then $G^{\bullet} \setminus \Sigma (T)$ is contained in the coset of a proper subgroup.

\item If $G^{\bullet}\setminus \Sigma (T)=\{g,h\}$ with $g \notin \langle h-g \rangle$, then $G^{\bullet}\setminus \Sigma (T)\subseteq g+\langle h-g \rangle$. 

\item If  $G^{\bullet}\setminus \Sigma (T)=\{g,h\}$ with $g, h  \in \langle h-g \rangle$, then there is no proper subgroup $H \subsetneq G$ and no $\alpha\in G\setminus H$  such that $G^{\bullet}\setminus \Sigma (T)=\{g,h\}\subseteq \alpha+H $.
\end{enumerate}  
\end{lemma}

\begin{equation} \label{def-nu(S)}
|T|\geq \nu(S)  \ \mbox{ implies } \ G^{\bullet} \setminus \Sigma (T) \subseteq \alpha+H \ \mbox {for some  subgroup $H  \subsetneq  G$ and some $\alpha\in G\setminus H$}.
\end{equation}

\begin{lemma} \label{6.2}
Let $G$ be a nontrivial, finite abelian group.
\begin{enumerate}
\item For every maximal zero-sum free sequence $S \in \mathcal F (G)$, we have $|S|-1 \le \nu (S) \le |S|$.

\item $\nu (G) = \max \{ \nu (S) \colon S \in \mathcal F (G) \ \text{is maximal zero-sum free} \}$.
\end{enumerate}
\end{lemma}

\begin{theorem} \label{6.5}
Let $G = C_2^2 \oplus C_6^2 \oplus C_{6n}$ with $n \in \N$, and let $(e_1, \ldots, e_5)$ be a basis of $G$ with $\ord (e_1) = \ord
(e_2) = 2$, $\ord (e_3) = \ord (e_4) = 6$ and $\ord (e_5) = 6n$.
\begin{enumerate}
\item The sequence 
      \begin{align*}
      S =   g_1 g_2 g_3 g_4 g_5 g_6  g_7^3 g_8^4 g_9^3 g_{10}^{6n-4} 
      \end{align*}
      is zero-sum free, where $g_1 = e_3+e_4+e_5$, $g_2 = e_2+e_5$, $g_3 =
      e_1+e_4$, $g_4 = e_2+e_4$, $g_5 = e_1+e_3$, $g_6=e_2+e_3$, $g_7 =
      e_1+e_2+2e_4+e_5$, $g_8 = 2e_3+e_4$, $g_9 = e_2+2e_3$ and $g_{10} =
      e_1+e_5$.  In particular, $\mathsf d (G) \ge |S|  = 6n+12 = \mathsf d^* (G)+1$.
      
\item $S$ is maximal zero-sum free.

\item $\nu (S) = |S| -1$.
\end{enumerate}      
\end{theorem}

\begin{proof}
1. Suppose that  $T = \prod_{i=1}^{10} g_i^{l_i}$ is a zero-sum
subsequence of $S$ with $l_i \in [0,1]$ for $i \in [1,6]$, $l_7, l_9
\in [0,3]$, $l_8 \in [0,4]$ and $l_{10} \in [0, 6n-4]$. We have to
show that $T = 1$. We start with the following system of
congruences:
\begin{align*}
\begin{aligned}
a_1 = l_3+l_5+l_7+l_{10} & \equiv 0 \mod 2 \\
a_2 = l_2+l_4+l_6+l_7+l_9 & \equiv 0 \mod 2 \\
a_3 = l_1+l_5+l_6+ 2l_8 + 2l_9 & \equiv 0 \mod 6 \\
a_4 = l_1+l_3+l_4+2l_7+l_8 & \equiv 0 \mod 6  \quad \text{and} \\
a_5 = l_1+l_2+l_7+l_{10} & \equiv 0 \mod 6n
\end{aligned}
\end{align*}
Note that $a_3 \in [0, 17]$, $a_4 \in [0, 13]$ and $a_5 \in [0,
6n+1]$. We distinguish the following two cases.


\smallskip
\noindent CASE 1: \,$a_5 = 0$.

Then $l_1=l_2=l_7=l_{10} = 0$ and $a_4 \in \{0, 6\}$.

Suppose that $a_4 = l_3+l_4+l_8=0$. Then $l_3=l_4=l_8=0$,
$a_1=l_5=0$, $a_2 = l_6+l_9 \equiv 0 \mod 2$ and $a_3 = l_6 + 2l_9
\equiv 0 \mod 6$. Thus $l_6 = 0$, $l_9 \equiv 0 \mod 2$ and $2 l_9
\equiv 0 \mod 6$. Therefore $l_9 = 0 = \sum_{i=1}^{10} l_i = |T|$.

Suppose that $a_4 = l_3+l_4+l_8=6$. Then $l_3=l_4=1$, $l_8 = 4$,
$l_5 = a_1-l_3=1$, $a_2 = 1 + l_6 + l_9 \equiv 0 \mod 2$ and $a_3 =
1 + l_6 + 8 + 2l_9 \equiv 0 \mod 6$. Then the last equation implies
 $l_6 = 1$, whence $a_2 \equiv l_9 \equiv 0 \mod 2$,  $4 + 2 l_9
\equiv 0 \mod 6$, and thus $l_9=4>3$, a contradiction.

\smallskip
\noindent CASE 2: \,$a_5 = 6n$.

Then $l_{10} \in \{6n-5, 6n-4\}$, and we distinguish two cases.

\smallskip
\noindent CASE 2.1: \,$l_{10} = 6n-5$.

Then $l_1=l_2=1$ and $l_7 = 3$. Then $a_4=1+l_3+l_4+6+l_8 \equiv 0
\mod 6$ and hence $l_3+l_4+l_8=5$.

If $l_8=3$, then $l_3=l_4=1$, $a_1=1+l_5+3+6n-5 \equiv 0 \mod 2$,
$l_5=1$, $a_2=1+1+l_6+3+l_9 \equiv 0 \mod 2$ and $a_3=1+1+l_6+6+2l_9
\equiv 0 \mod 6$. The last equation implies $l_6=0$, hence $1+l_9
\equiv 0 \mod 2$ and $2+2l_9 \equiv 0 \mod 6$, a contradiction.

Suppose that $l_8=4$. If $l_4=0$, then $l_3=1$, $a_1=l_3+l_5 \equiv
0 \mod 2$, $l_5=1$, $a_2=1+0+l_6+3+l_9 \equiv 0 \mod 2$,
$a_3=1+1+l_6+8+2l_9 \equiv 0 \mod 6$, hence $2 \t l_9$, $l_6=0$ and
$10+2l_9 \equiv 0 \mod 6$, a contradiction. If $l_4=1$, then
$l_3=0$, $a_1=l_3+l_5 \equiv 0 \mod 2$, $l_5=0$, $a_2=1+1+l_6+3+l_9
\equiv 0 \mod 2$, $a_3=1+0+l_6+8+2l_9 \equiv 0 \mod 6$, hence
$l_6=1$, $2 \t l_9$ and $10+2l_9 \equiv 0 \mod 6$, a contradiction.

\smallskip
\noindent CASE 2.2: \,$l_{10} = 6n-4$.

Then $l_7 \in \{2,3\}$.

\smallskip
\noindent CASE 2.2.1: \,$l_{7} = 2$.

Then $l_1=l_2=1$, $2 \t a_1$, $l_3=l_5$, $a_4 = 1+l_3+l_4+4+l_8
\equiv 0 \mod 6$ and hence $l_3+l_4+l_8=1$.

If $l_3=1$, then $l_4=l_8=0$, $a_2=1+l_6+2+l_9 \equiv 0 \mod 2$,
$a_3=1+1+l_6+2l_9 \equiv 0 \mod 6$, $l_6=0$, $l_9$ odd and $2+2l_9
\equiv 0 \mod 6$, a contradiction.

If $l_4=1$, then $l_3=l_8=0$, $a_2=1+1+l_6+2+l_9 \equiv 0 \mod 2$,
$a_3 = 1+l_6+2l_9 \equiv 0 \mod 6$, $l_6=1$, $l_9$ odd and $2+2l_9
\equiv 0 \mod 6$, a contradiction.

If $l_8=1$, then $l_3=l_4=0$, $a_2=1+l_6+2+l_9 \equiv 0 \mod 2$,
$a_3=1+l_6+2+2l_9 \equiv 0 \mod 6$, $l_6=1$, $l_9$ even and $4+2l_9
\equiv 0 \mod 6$, a contradiction.

\smallskip
\noindent CASE 2.2.2: \,$l_{7} = 3$.

Then $l_1+l_2 = 1$, $a_1 \equiv l_3+l_5 \equiv 1 \mod 2$ (hence
$(l_1,l_2,l_3,l_5) \in \{(1,0,1,0), (1,0,0,1), (0,1,1,0),$ $
(0,1,0,1)\}$), and $a_2 \equiv l_2+l_4+l_6+l_9 \equiv 1 \mod 2$,
$a_3 \equiv l_1+l_5+l_6+2l_8+2l_9 \equiv 0 \mod 6$ and $a_4 \equiv
l_1+l_3+l_4+l_8 \equiv 0 \mod 6$.

If $(l_1,l_2,l_3,l_5) = (1,0,1,0)$, then $l_4+l_6+l_9 \equiv 1 \mod
2$, $1+l_6+2l_8+2l_9 \equiv 0 \mod 6$ and $2+l_4+l_8 \equiv 0 \mod
6$, hence $l_6 = 1$, $2 \t l_4+l_9$, $3 \t 1+l_8+l_9$ and $6 \t
2+l_4+l_8$, a contradiction.

If $(l_1,l_2,l_3,l_5) = (1,0,0,1)$, then $l_4+l_6+l_9 \equiv 1 \mod
2$, $2+l_6+2l_8+2l_9 \equiv 0 \mod 6$ and $1+l_4+l_8 \equiv 0 \mod
6$, hence $l_6 = 0$, $l_4=1$ and $l_8=4$. The first congruence forces $l_9$ to be even, whereas the second gives $l_9\equiv1\mod3$; this is impossible for $l_9\in[0,3]$. 

If $(l_1,l_2,l_3,l_5) = (0,1,1,0)$, then $l_4+l_6+l_9 \equiv 0 \mod
2$, $l_6+2l_8+2l_9 \equiv 0 \mod 6$ and $1+l_4+l_8 \equiv 0 \mod 6$,
hence $l_6=0$, $2 \t l_4+l_9$, $3 \t l_8+l_9$ and $6 \t 1+l_4+l_8$;
so $l_4=1$, $l_8=4$ and $l_9=2$, a contradiction.

If $(l_1,l_2,l_3,l_5) = (0,1,0,1)$, then $l_4+l_6+l_9 \equiv 0 \mod
2$, $1+l_6+2l_8+2l_9 \equiv 0 \mod 6$ and $l_4+l_8 \equiv 0 \mod 6$,
hence $l_6=1$, $2 \t 1+l_4+l_9$, $3 \t 1+l_8+l_9$ and $6 \t
l_4+l_8$, a contradiction.

\smallskip
2. In order to prove that $S$ is maximal zero-sum free 
we  introduce a further basis for $G$. We define
$e_5' = e_1 + e_5$ and observe that 
$(e_1,e_2,e_3,e_4,e_5')$ is  a basis of $G$. Then we have
\begin{align*}
	g_1&=e_1+e_3+e_4+e_5',\quad g_2=e_1+e_2+e_5',\quad g_3=e_1+e_4,\quad g_4=e_2+e_4,\\
	g_5&=e_1+e_3,\quad g_6=e_2+e_3,\quad g_7=e_2+2e_4+e_5',\\
	g_8&=2e_3+e_4,\quad g_9=e_2+2e_3,\quad g_{10}=e_5'.
\end{align*}
For any element $x_1e_1+x_2e_2+x_3e_3+x_4e_4+x_5e_5' \in G$, we use the compact coordinate notation
\begin{align*}
	(x_1,x_2,x_3,x_4,x_5)^t &:= x_1e_1+x_2e_2+x_3e_3+x_4e_4+x_5e_5'\\
	&=(e_1,e_2,e_3,e_4,e_5')\cdot (x_1,x_2,x_3,x_4,x_5)^t.
\end{align*}
We set
\[
S = S'  (e_5')^{6n-4} \quad \text{with} \quad S' = g_1 g_2 g_3 g_4 g_5 g_6 \, g_7^3 g_8^4 g_9^3 \,.
\]
We have 
\[
\Sigma^*\bigl((e_5')^{6n-4}\bigr) = \langle e_5' \rangle \setminus \{-e_5',-2e_5',-3e_5'\} \,,
\]
and for 
$T = g_3 g_4 g_5 g_6 \, g_8^4 g_9^3$ we have
\begin{align*}
	\Sigma^*(T)=
	\{ \left(\begin{array}{l|l|l|l|l|l|l|l|l|l|l|l|l|l}
		0 & 1 & 0 & 1 & 0 & 1 & 0 & 1 & 1 & 0 & 1 & 0 & 1 & 0\\
		0 & 0 & 1 & 0 & 1 & 1 & 0 & 1 & 1 & 1 & 0 & 1 & 0 & 0\\
		0 & 0 & 0 & 1 & 1 & 0 & 1 & 1 & 2 & 1 & 1 & 2 & 2 & 2\\
		0 & 1 & 1 & 0 & 0 & 2 & 1 & 1 & 0 & 2 & 2 & 1 & 1 & 2\\
		0 & 0 & 0 & 0 & 0 & 0 & 0 & 0 & 0 & 0 & 0 & 0 & 0 & 0
	\end{array}\right)
	\} \\
	+  \{
	\left(\begin{array}{l|l|l|l|l}
		0 & 0 & 0 & 0 & 0\\
		0 & 0 & 0 & 0 & 0\\
		0 & 2 & 4 & 0 & 2\\
		0 & 1 & 2 & 3 & 4\\
		0 & 0 & 0 & 0 & 0
	\end{array}\right)
	\}+  \{
	\left(\begin{array}{l|l|l|l}
		0 & 0 & 0 & 0\\
		0 & 1 & 0 & 1\\
		0 & 2 & 4 & 0\\
		0 & 0 & 0 & 0\\
		0 & 0 & 0 & 0
	\end{array}\right)
	\}.
\end{align*}
%The complement of $\Sigma(T)$ in $\langle e_1,e_2,e_3,e_4 \rangle \setminus \{0\}$ satisfies
Therefore, $ \langle e_1, \ldots, e_4 \rangle^{\bullet} \setminus \Sigma (T)=\langle e_1, \ldots, e_4 \rangle \setminus \Sigma^* (T)$ is a subset of 
\begin{align*}
\{\left(\begin{array}{l|l|l|l|l|l}
	1 & 0 & 1 & 0 & 1 & 1\\
	0 & 1 & 0 & 0 & 1 & 1\\
	0 & 5 & 5 & 5 & 5 & 4\\
	0 & 1 & 1 & 5 & 5 & 0\\
	0 & 0 & 0 & 0 & 0 & 0
\end{array}\right)
\} \subseteq 
\left(\begin{array}{l}
	1 \\ 0 \\ 0 \\ 0 \\ 0 
\end{array}\right)
+\left\langle \left(\begin{array}{l}
	1 \\ 1 \\ 0 \\ 0 \\ 0 
\end{array}\right),
\left(\begin{array}{l}
	1 \\ 0 \\ 2 \\ 0 \\ 0 
\end{array}\right),
\left(\begin{array}{l}
	0 \\ 0 \\ 5 \\ 1 \\ 0 
\end{array}\right) \right\rangle.
\end{align*}


Let $G_1\cong C_2^2\oplus C_6^3$ have a basis $(e_1,e_2,e_3,e_4,e_0)$ whose elements have orders $2,2,6,6,6$, respectively.
We define a group homomorphism
\begin{align*}
f: G \to G_1,\quad x_1e_1+x_2e_2+x_3e_3+x_4e_4+x_5e_5' \longmapsto x_1e_1+x_2e_2+x_3e_3+x_4e_4+x_5 e_0.
\end{align*}


Since $\Sigma(S')\subseteq \{0,e_5',...,5e_5'\}+ \langle e_1,...,e_4\rangle$, the restriction of $f$ to $\Sigma(S')$ is injective, and 
\[
f\left( (e_1,e_2,e_3,e_4,e_5')\cdot (x_1,x_2,x_3,x_4,x_5)^t\right) =(e_1,e_2,e_3,e_4,e_0)\cdot (x_1,x_2,x_3,x_4,x_5)^t
\]
with $x_5\in [0,5]$. For the subset
$$C = (\{0,e_5',...,5e_5'\}+ \langle e_1,...,e_4\rangle)^{\bullet} \setminus \Sigma(S') \ \subseteq \ G^{\bullet}$$
we obtain that
\begin{align*}
f(C)=
\{(e_1,e_2,e_3,e_4,e_0)\cdot 
\left(\begin{array}{l|l|l|l|l|l|l|l|l|l|l|l|l|l}
	0 & 0 & 0 & 1 & 0 & 1 & 0 & 0 & 1 & 1 & 1 & 1 & 1 & 1\\
	0 & 0 & 0 & 1 & 1 & 0 & 1 & 1 & 0 & 0 & 1 & 1 & 1 & 0\\
	0 & 0 & 5 & 5 & 0 & 0 & 5 & 5 & 1 & 5 & 0 & 4 & 5 & 0\\
	0 & 0 & 5 & 5 & 2 & 2 & 1 & 1 & 1 & 1 & 0 & 0 & 1 & 0\\
	4 & 5 & 0 & 0 & 5 & 5 & 0 & 1 & 5 & 0 & 5 & 0 & 5 & 0
\end{array}\right)
\} = G_1^{\bullet} \setminus \Sigma (f(S')) \,,
\end{align*}

whence

\begin{align*}
C=
\{(e_1,e_2,e_3,e_4,e_5')\cdot 
\left(\begin{array}{l|l|l|l|l|l|l|l|l|l|l|l|l|l}
	0 & 0 & 0 & 1 & 0 & 1 & 0 & 0 & 1 & 1 & 1 & 1 & 1 & 1\\
	0 & 0 & 0 & 1 & 1 & 0 & 1 & 1 & 0 & 0 & 1 & 1 & 1 & 0\\
	0 & 0 & 5 & 5 & 0 & 0 & 5 & 5 & 1 & 5 & 0 & 4 & 5 & 0\\
	0 & 0 & 5 & 5 & 2 & 2 & 1 & 1 & 1 & 1 & 0 & 0 & 1 & 0\\
	4 & 5 & 0 & 0 & 5 & 5 & 0 & 1 & 5 & 0 & 5 & 0 & 5 & 0
\end{array}\right)
\},
\end{align*}


We now claim that $f$ induces an injective map from the set 
\begin{align*}
A := G^\bullet \setminus \Sigma\left(S' g_{10}^{6n-6}\right)
\end{align*}
to the set
\begin{align*}
B := G_1^\bullet \setminus \Sigma\left(f(S')\right) = f(C).
\end{align*}
We  recall that $g_{10} = e_5'$, whence $\Sigma^* \big( g_{10}^{6n-6} \big) = \{0, e_5', 2e_5', \dots, (6n-6)e_5'\}$, and we continue with the following two assertions.
\begin{enumerate}
\item[{\bf A1.}\,] $f (A) \subseteq B$.

\item[{\bf A2.}\,] $f$ is injective on $A$.
\end{enumerate}

\smallskip

{\it Proof of \,{\bf A1}}. Let $g = (x_1, x_2, x_3, x_4, x_5)^t \in A$. Since
	$g \notin \Sigma(S' g_{10}^{6n-6})$, we claim that
	\begin{equation}\label{eq:not-in-SigmaSprime}
	g + m e_5' \notin \Sigma(S')
	\quad \text{for all } m \in [6, 6n].
	\end{equation}
	Indeed, if $g + m e_5' \in \Sigma(S')$ for some $m \in [6, 6n]$, then
	$-m e_5' \in \Sigma^*(g_{10}^{6n-6})$ (since
	$-m \equiv 6n - m \in [0, 6n-6] \mod{6n}$), and rearranging gives
	$g = (g + m e_5') + (-m e_5') \in \Sigma(S' g_{10}^{6n-6})$,
	contradicting $g \in A$. In particular, taking $m = 6k$ for
	$k \in [1,n]$, we obtain $g + 6k e_5' \notin \Sigma(S')$.
	
	Let $s = x_5 \bmod 6$, so that $6 \mid (6n + s - x_5)$ and
	$6n + s - x_5 \in [1, 6n]$. Set $k_0 = (6n + s - x_5)/6$. Then
	$g + 6k_0 e_5'$ has $e_5'$-coordinate $s \in [0, 5]$, so
	$g + 6k_0 e_5' \in
	\bigl(\{0, e_5', \dots, 5e_5'\} + \langle e_1, \dots, e_4\rangle
	\bigr)^{\bullet}.$
	
	On the other hand, $g + 6k_0 e_5' \notin \Sigma(S')$,
	hence $g + 6k_0 e_5' \in C$. Since $f(e_5') = e_0$ and
	$\ord(e_0) = 6$, we have
	$f(g) = f(g) + k_0 (6 e_0) = f(g + 6k_0 e_5') \in f(C) = B$,
	as required.
	
{\it Proof of \,{\bf A2}}.	 Assume, to the contrary, that there exist distinct $g, h \in A$ with
	$f(g) = f(h)$. Then $h - g \in \ker(f) = \langle 6 e_5' \rangle$, so
	$h = g + 6k_1 e_5'$ for some integer $k_1$ with $1 \leq k_1 \leq n-1$.
	
	Applying~\eqref{eq:not-in-SigmaSprime} to $g$, we have
	$g + j e_5' \notin \Sigma(S')$ for every
	$j \in [6,6n]$.
	Applying~\eqref{eq:not-in-SigmaSprime} to $h$ and using
	$h + m e_5' = g + (6k_1 + m) e_5'$ yields
	$g + j e_5' \notin \Sigma(S')$ for every
	$j \in \mathbb{Z}/6n\mathbb{Z} \setminus \{6k_1 + 1, \dots, 6k_1 + 5\}$.
	Since $1 \leq k_1 \leq n-1$, we have $6k_1 + 1 \geq 7$ and
	$6k_1 + 5 \leq 6n - 1$, so the two exempt sets $\{1, \dots, 5\}$ and
	$\{6k_1 + 1, \dots, 6k_1 + 5\}$ are disjoint in $\mathbb{Z}/6n\mathbb{Z}$.
	Combining the two conditions, $g + j e_5' \notin \Sigma(S')$ for
	\emph{every} $j \in \mathbb{Z}/6n\mathbb{Z}$, that is, $(g + \langle e_5' \rangle) \cap \Sigma(S') = \emptyset.$
	
	In particular, the element
	$g' := (x_1, x_2, x_3, x_4, 2)^t = g + (2 - x_5) e_5'$
	lies in $g + \langle e_5'\rangle$ and satisfies $g' \notin \Sigma(S')$.
	Since its $e_5'$-coordinate is $2 \in [0, 5]$, we have
	$g' \in C$. However, inspection of the explicit description
	of $C$ above shows that the $e_5'$-coordinates of its
	elements (the last row of the matrix) take values only in
	$\{0, 1, 4, 5\}$, never $2$. This contradicts $g' \in C$.
	Therefore $f$ is injective on $A$.


\medskip
Combining {\bf A1} and {\bf A2} with
$
\Sigma^*\bigl((e_5')^{6n-6}\bigr)
=\{0,e_5',\ldots,(6n-6)e_5'\},
$
and translating in the $e_5'$-coordinate, we obtain the following
list of the only possible lifts of the elements of set $B$.
\begin{align*}
G^{\bullet} \setminus \Sigma (S'g_{10}^{6n-6}) \subseteq
\{\left(\begin{array}{l|l|l|l|l|l|l|l|l|l|l|l|l|l}
	0 & 0 & 0 & 1 & 0 & 1 & 0 & 0 & 1 & 1 & 1 & 1 & 1 & 1\\
	0 & 0 & 0 & 1 & 1 & 0 & 1 & 1 & 0 & 0 & 1 & 1 & 1 & 0\\
	0 & 0 & 5 & 5 & 0 & 0 & 5 & 5 & 1 & 5 & 0 & 4 & 5 & 0\\
	0 & 0 & 5 & 5 & 2 & 2 & 1 & 1 & 1 & 1 & 0 & 0 & 1 & 0\\
	6n-2 & 6n-1 & 0 & 0 & 6n-1 & 6n-1 & 0 & 1 & 6n-1 & 0 & 6n-1 & 0 & 6n-1 & 0
\end{array}\right)
\} \,,
\end{align*}
whence $G^{\bullet} \setminus \Sigma (S'g_{10}^{6n-5})\subseteq \{(0,0,0,0,6n-1),(0,1,5,1,1)\}$, and thus $\Sigma(S)=G^{\bullet}$. 

\bigskip
3. By Lemma \ref{6.2}, it suffices to verify that  $\nu (S) \le |S|-1$. In order to do so, we prove that
$|S| - 1$ satisfies \eqref{def-nu(S)}. Let $T \in \mathcal F (G)$ with $|T|=|S|-1$, say $S = gT$ for some $g \in \supp (S)$, and $g\notin 2G$.  Note that $\ord(g) \in \{6, 6n\}$  and  $\sigma(S) = e_2 + 5e_3 + e_4+e_5'$ has order $6n$. We distinguish two cases.

\smallskip
\noindent
CASE 1: $g \neq g_1$.

Using the same reduction of the $e_5'$-coordinate modulo
$\langle 6e_5'\rangle$ as in Part 2, it suffices to carry out the
verification in $
G/\langle 6e_5'\rangle\cong C_2^2\oplus C_6^3.$
A direct computation in this quotient, together with the injective lifting argument in Part 2 applied to $Sg^{-1}$, shows that 
\[
	G^{\bullet} \setminus \Sigma(Sg^{-1})=\{-g, \sigma (S) \} \,.
\]
If $\ord(g)=6n$, then $g\in\{g_2,g_7,g_{10} \}$ whence $g \in e_5'+\langle e_1,...,e_4\rangle$. Thus $\sigma (S) \notin \langle g+ \sigma (S) \rangle \subseteq \langle e_1,...,e_4,2e_5'\rangle$, whence $\{-g, \sigma (S) \} \subseteq e_5' + \langle e_1,...,e_4,2e_5'\rangle$.
		
If $\ord(g)=6$, then $g+\sigma (S) \in e_5'+\langle e_1,...,e_4\rangle$ and $\ord(g+ \sigma (S))=6n$. Since $\langle e_1,...,e_4\rangle \cap \langle g+\sigma (S) \rangle=\{0\}$ and $g\ne 0$, it follows that  $-g\notin \langle g+\sigma (S) \rangle$. Thus Lemma \ref{2.3}(3) implies the assertion similarly.


	
\smallskip
\noindent
CASE 2: $g=g_1$. 

The proof of Part 2 (or direct computation) shows that
	\begin{align*}
	\begin{aligned}
	G^{\bullet} \setminus \Sigma(Sg_1^{-1})
	&=
	\bigl\{
	e_1+5e_3+5e_4+(6n-1)e_5',\, e_1+5e_3+e_4+e_5',\\
	&\qquad e_2+5e_3+5e_4+(6n-1)e_5',\, e_2+5e_3+e_4+e_5'
	\bigr\}.
	\end{aligned}
	\end{align*}
Thus we obtain that
\[
	G^{\bullet} \setminus \Sigma(Sg_1^{-1})\subseteq
	g_1 + \langle e_1,e_2,e_3, e_4,2e_5' \rangle \,. \qedhere
\]
\end{proof}

\end{document}
